% C-pass-interactions.tex — Annex C: Pass Interactions
\chapter{Pass Interactions}
\label{annex:C-pass-interactions}
\index{pass interactions}
\index{semantic analysis!passes}

\pnum
The Lain compiler's semantic analysis is not a sequence of fully independent
passes that each transform the AST and hand it off to the next.  Instead it is
two \emph{phases} — resolve and walk — orchestrated by a single function
\cname{sema\_resolve\_module} in \srcref{src/sema.h}, with tightly coupled
sharing of mutable global state between the analysis components within each
phase.  This annex describes how those components interact with one another and
with the shared state, and how the two-phase structure enforces the dependency
ordering between analyses.

% -------------------------------------------------------------------------
\section{Two-phase structure}
\label{sec:ci-twophase}
\index{sema\_walk\_phase}
% -------------------------------------------------------------------------

\pnum
The boolean flag \cname{sema\_walk\_phase} (declared in \srcref{src/sema.h})
partitions all semantic analysis into two epochs per function body:

\begin{enumerate}
\item \textbf{Resolve phase} (\cname{sema\_walk\_phase = false}).
  \cname{sema\_resolve\_stmt} is called on every statement.
  Name resolution, UFCS desugaring, generic instantiation, \cname{comptime if}
  branch selection, and the initial call to \cname{sema\_infer\_expr} all occur
  here.  VRA range propagation, bounds checks, and linearity checks are
  \emph{not} performed.

\item \textbf{Walk phase} (\cname{sema\_walk\_phase = true}).
  \cname{walk\_stmt} is called on every statement a second time.  This invokes
  \cname{sema\_infer\_expr} again, which now sees all \cname{->type} annotations
  established during resolve, and activates VRA range propagation, in-guard
  tracking, bounds checking, and the termination-measure verifier.
\end{enumerate}

\pnum
The guard idiom used by walk-only checks is:

\begin{ccode}
if (sema_walk_phase) {
    /* bounds check, VRA check, etc. */
}
\end{ccode}

\noindent
This prevents false positives during resolve, when type information may still
be incomplete (e.g.\ a generic is being instantiated, or a forward reference
is being resolved).

% -------------------------------------------------------------------------
\section{Shared global state}
\label{sec:ci-globals}
\index{global state!semantic}
% -------------------------------------------------------------------------

\pnum
All semantic passes share the following global tables, declared in
\srcref{src/sema.h} before the sub-header includes:

\begin{center}
\begin{tabular}{lll}
\toprule
\textbf{Global} & \textbf{Type} & \textbf{Consumers} \\
\midrule
\cname{sema\_globals}  & \cname{HashTable} (4096 buckets) & Resolve, typecheck \\
\cname{sema\_locals}   & \cname{HashTable} (4096 buckets) & Resolve, typecheck \\
\cname{sema\_scope}    & \cname{ScopeFrame}               & Resolve, linearity \\
\cname{sema\_ranges}   & \cname{RangeTable *}             & VRA, bounds \\
\cname{sema\_in\_guards} & \cname{InGuardEntry *}         & VRA, bounds \\
\cname{sema\_ltable}   & \cname{LTable *}                 & Linearity \\
\cname{sema\_btable}   & \cname{BorrowTable *}            & Borrows \\
\cname{sema\_source\_text} & \cname{const char *}         & Diagnostics \\
\cname{sema\_source\_file} & \cname{const char *}         & Diagnostics \\
\bottomrule
\end{tabular}
\end{center}

\pnum
The symbol tables (\cname{sema\_globals}, \cname{sema\_locals}) are written by
the resolver (name binding and UFCS desugaring) and read by the typechecker
(callee lookup, struct field lookup).  The range table (\cname{sema\_ranges})
is written by VRA and read by the bounds checker.  The in-guard table
(\cname{sema\_in\_guards}) is written by \cname{sema\_push\_in\_guards} and
read by the bounds checker.  The linearity table (\cname{sema\_ltable}) and
borrow table (\cname{sema\_btable}) are written and read exclusively by their
respective passes.

% -------------------------------------------------------------------------
\section{Resolve phase interactions}
\label{sec:ci-resolve}
\index{resolve phase}
% -------------------------------------------------------------------------

\pnum
Within the resolve phase the call graph is:

\begin{enumerate}
\item \cname{sema\_resolve\_stmt} dispatches on \cname{StmtKind}.  For every
  expression inside the statement it calls \cname{sema\_resolve\_expr}, which
  handles UFCS transformation and generic call expansion.

\item \cname{sema\_resolve\_expr} calls \cname{sema\_infer\_expr} on any
  sub-expression whose type is needed for resolution decisions (e.g.\ the
  receiver type of a member call is needed to determine the UFCS candidate).
  This creates a partial re-entrant dependency between resolve and typecheck.

\item For generic calls (\cname{EXPR\_CALL} on a comptime-parameterised
  \cname{DECL\_FUNCTION}), \cname{sema\_resolve\_expr} calls into
  \srcref{src/sema/comptime.h} to evaluate comptime arguments, then calls
  \srcref{src/sema/generic.h} to clone and substitute the function body, then
  recursively calls \cname{sema\_resolve\_module} on the specialisation.  This
  means generic instantiation is fully recursive and nested within resolve.

\item For \cname{STMT\_COMPTIME\_IF}, \cname{sema\_resolve\_stmt} evaluates
  the condition via \cname{comptime\_evaluate\_expr}, sets
  \cname{is\_taken}, and recursively resolves only the selected branch.
  The unselected branch is never visited.
\end{enumerate}

% -------------------------------------------------------------------------
\section{Walk phase interactions}
\label{sec:ci-walk}
\index{walk phase}
% -------------------------------------------------------------------------

\pnum
The walk phase processes each statement in the same statement-list order as
resolve.  \cname{walk\_stmt} calls \cname{sema\_infer\_expr}, which in turn
calls the following analyses:

\begin{center}
\begin{tabular}{lll}
\toprule
\textbf{Analysis} & \textbf{Entry point} & \textbf{Called from} \\
\midrule
Type inference      & \cname{sema\_infer\_expr}     & \cname{walk\_stmt} \\
VRA propagation     & \cname{range\_set}/\cname{range\_eval} & Inside \cname{sema\_infer\_expr} \\
In-guard push       & \cname{sema\_push\_in\_guards} & \cname{walk\_stmt} (STMT\_WHILE/STMT\_IF) \\
Bounds check        & \cname{sema\_check\_bounds}    & Inside \cname{sema\_infer\_expr} (EXPR\_INDEX) \\
Linearity check     & \cname{linearity\_check\_stmt} & \cname{walk\_stmt} (after infer) \\
Borrow check        & \cname{borrow\_register}       & Inside \cname{linearity\_check\_stmt} \\
Termination verify  & bounded-while verifier         & \cname{walk\_stmt} (STMT\_WHILE with measure) \\
\bottomrule
\end{tabular}
\end{center}

\pnum
The ordering within a single walk-phase statement is significant:
\begin{enumerate}
\item \cname{sema\_infer\_expr} runs first.  It populates \cname{->type} on
  the expression tree and propagates range information into \cname{sema\_ranges}.
  Bounds checks are triggered as a side-effect of inferring \cname{EXPR\_INDEX}
  expressions.
\item In-guards are pushed immediately before the loop/if body is walked.
  The guard for a condition \lain{idx in arr} is added to \cname{sema\_in\_guards}
  so that bounds checks on \lain{arr[idx]} inside the body see the proof.
\item Linearity checking (\cname{linearity\_check\_stmt}) runs after type
  inference for the same statement, so it sees the \cname{->type} annotations
  needed to determine linearity.  Borrow registration is a sub-step of linearity
  checking: shared borrows are registered when a variable is passed to a shared
  parameter, mutable borrows when passed to a \lain{var} parameter.
\item The bounded-while verifier is invoked at \cname{STMT\_WHILE} nodes that
  carry a non-null \cname{measure} field.  It reads the
  current \cname{sema\_ranges} state to evaluate non-negativity and reads
  the condition subtree to extract in-guard patterns for the decrease proof.
\end{enumerate}

% -------------------------------------------------------------------------
\section{VRA and bounds interaction}
\label{sec:ci-vra-bounds}
\index{VRA!interaction with bounds}
% -------------------------------------------------------------------------

\pnum
VRA (\cref{chap:14-vra}) and the bounds checker (\cref{chap:15-bounds}) share
the \cname{sema\_ranges} table and the \cname{sema\_in\_guards} linked list.
Their interaction is:

\begin{enumerate}
\item During \cname{sema\_infer\_expr} for \cname{EXPR\_BINARY} assignment,
  \cname{range\_set} updates the range of the target variable.
\item When \cname{sema\_infer\_expr} encounters \cname{EXPR\_INDEX}, it reads
  the range of the index expression via \cname{range\_eval} and calls
  \cname{sema\_check\_bounds}.
\item \cname{sema\_check\_bounds} first checks the in-guard table for a proven
  \texttt{(index, container)} pair.  If found, it returns without emitting a
  runtime check.  If not, it queries \cname{sema\_ranges} for the range.
  Static ranges trigger either a compile-time error or a silent pass;
  unknown ranges trigger emission of a C runtime bounds-check expression.
\end{enumerate}

\pnum
This means VRA must run before the bounds checker for any given index
expression, and the in-guard table must be populated before the loop body is
walked.  Both of these orderings are guaranteed by the sequential evaluation
inside \cname{sema\_infer\_expr} and the explicit push-before-body structure
in \cname{walk\_stmt}.

% -------------------------------------------------------------------------
\section{Linearity and borrows interaction}
\label{sec:ci-lin-borrow}
\index{linearity!interaction with borrows}
% -------------------------------------------------------------------------

\pnum
The linearity checker (\cref{chap:12-linearity}) and the borrow checker
(\cref{chap:13-borrows}) share \cname{sema\_ltable} and \cname{sema\_btable}
but do not call each other directly.  Their interaction is through the following
read/write protocol:

\begin{itemize}
\item When the linearity checker encounters a move (\cname{EXPR\_MOVE} or an
  implicit transfer into a \cname{MODE\_OWNED} parameter), it calls
  \cname{borrow\_is\_live(sema\_btable, id)} to verify no borrow is active
  before marking the variable consumed.  If a borrow is live, \diagcode{E008}
  is emitted.

\item When a variable is passed to a shared (\cname{MODE\_SHARED}) parameter,
  the linearity checker calls \cname{borrow\_register} with \cname{MODE\_SHARED}
  and registers the two-phase RESERVED borrow.  After all arguments are
  evaluated, \cname{borrow\_promote\_to\_active} upgrades the borrow to ACTIVE.

\item When a variable is passed to a mutable (\cname{MODE\_MUTABLE}) parameter,
  the linearity checker calls \cname{borrow\_register} with \cname{MODE\_MUTABLE}.
  The borrow checker verifies no conflicting borrow exists (\diagcode{E004}) before
  registering.

\item At scope exit, the linearity checker scans \cname{sema\_ltable} for
  unconsumed linear variables (\diagcode{E003}), and the borrow checker
  automatically expires borrows whose region no longer contains the current
  statement index.
\end{itemize}

% -------------------------------------------------------------------------
\section{Mutual recursion detection}
\label{sec:ci-mrec}
\index{mutual recursion!detection}
% -------------------------------------------------------------------------

\pnum
Mutual recursion detection (the DFS described in \cref{chap:09-sema-orchestration})
runs as a standalone pre-pass over the resolved declaration list, after both
phases have completed.  It does not share state with the other analyses except
reading the resolved \cname{->decl} pointers set during resolve.  Its output is
purely diagnostic: \diagcode{E011} is emitted and the exit code is set if a cycle
involving a \lain{func} is found.

% -------------------------------------------------------------------------
\section{Pass dependency DAG}
\label{sec:ci-dag}
\index{pass dependency}
% -------------------------------------------------------------------------

\pnum
The dependency relationships between passes can be summarised as a directed
acyclic graph:

\begin{center}
\begin{tikzpicture}[
  node distance=1.5cm and 2.5cm,
  every node/.style={font=\small},
  box/.style={draw, rounded corners=2pt, minimum width=2.8cm, minimum height=0.7cm,
              fill=white, align=center},
  arrow/.style={->, >=stealth, thin}
]
\node[box] (lex)    {Lexer};
\node[box, right=of lex] (parse) {Parser};
\node[box, right=of parse] (res)   {Resolve};
\node[box, below=of res]   (infer) {Type inference};
\node[box, left=of infer]  (vra)   {VRA};
\node[box, left=of vra]    (bounds){Bounds};
\node[box, below=of infer] (lin)   {Linearity};
\node[box, right=of lin]   (borrow){Borrows};
\node[box, below=of vra]   (term)  {Termination};
\node[box, below=of bounds] (mrec) {Mutual rec.};
\node[box, right=of borrow] (emit) {Emit};

\draw[arrow] (lex)    -- (parse);
\draw[arrow] (parse)  -- (res);
\draw[arrow] (res)    -- (infer);
\draw[arrow] (infer)  -- (vra);
\draw[arrow] (vra)    -- (bounds);
\draw[arrow] (infer)  -- (lin);
\draw[arrow] (lin)    -- (borrow);
\draw[arrow] (vra)    -- (term);
\draw[arrow] (res)    -- (mrec);
\draw[arrow] (lin)    -- (emit);
\draw[arrow] (borrow) -- (emit);
\end{tikzpicture}
\end{center}

\pnum
Edges represent ``must run before'' ordering.  All arcs involving the
walk-phase analyses (VRA, bounds, linearity, borrows, termination) are
implicitly conditioned on the walk phase having begun
(\cname{sema\_walk\_phase = true}).
